\subsection{Роль операционной системы при выполнении параллельной программы}
\label{subsec:os-role}

Параллельная программа взаимодействует с двумя планировщиками. Среда выполнения распределяет
  программные задачи между потоками ОС, а операционная система назначает эти потоки доступным
  ядрам. Поэтому время выполнения зависит от числа одновременно готовых задач, способа их
  планирования и состояния системы.

Потоки одного процесса используют общее адресное пространство, но у каждого есть собственный стек
  и состояние исполнения. Поток, ожидающий операцию ввода-вывода или блокировку, временно не
  занимает процессор; его место может получить другой готовый поток. Когда готовых потоков больше,
  чем аппаратных потоков CPU, операционная система распределяет между ними процессорное
  время~\cite{apple_threading_guide}.

\subsection{Планирование потоков в Linux и macOS}
\label{subsec:os-scheduling}

В Linux поведение планировщика зависит от версии ядра и класса планирования. Например, начиная
  с Linux 6.6 основной справедливый планировщик постепенно переходит к алгоритму
  EEVDF~\cite{linux_eevdf}. Это показывает, почему результаты одной конфигурации Linux нельзя
  автоматически переносить на все версии системы.

В macOS низкоуровневую основу потоков составляют потоки Mach. Программа также может использовать
  интерфейс POSIX threads, а библиотеки Apple предлагают очереди задач более высокого
  уровня~\cite{apple_threading_guide}.
  
% Технологии Cocoa и Grand Central Dispatch здесь не используются: обе сравниваемые реализации написаны на Go.

На распределение процессорного времени влияют также фоновые процессы, приоритеты и режим питания. Из-за
  этого результаты отдельных запусков могут различаться даже на одном компьютере. Повторные
  измерения помогают отделить устойчивую разницу между вариантами программы от случайных колебаний.

% TODO: добавить схему уровней планирования: задачи Go -> потоки ОС -> ядра/аппаратные потоки CPU.

\subsection{Среда выполнения Go и отображение горутин на потоки ОС}
\label{subsec:go-runtime}

Планировщик Go описывает выполнение через три сущности: G — горутина, M — поток ОС,
  P — ресурс, необходимый потоку M для выполнения кода Go. Число P задаётся параметром
  \texttt{GOMAXPROCS}. Среда выполнения распределяет готовые горутины между потоками ОС; поэтому
  множество горутин не требует такого же числа одновременно работающих системных
  потоков~\cite{go_runtime_hacking,go_runtime_scheduler}.

Если системный поток блокируется в системном вызове, его P может перейти другому потоку, чтобы
  готовые горутины продолжили выполнение. При этом общее число созданных потоков ОС не обязано
  совпадать с \texttt{GOMAXPROCS}. Такое отображение горутин на потоки соответствует модели M:N
  в рамках данного эксперимента~\cite{go_runtime_hacking}.

Для приближения модели 1:1 используется \texttt{runtime.LockOSThread}. Эта функция закрепляет
  горутину за текущим потоком ОС; до снятия закрепления другие горутины не выполняются на этом
  потоке~\cite{go_runtime_package}. Если горутины, выполняющие части вычисления, закреплены за
  отдельными потоками ОС, каждой такой горутине соответствует свой системный поток. Управление
  потоками остаётся частью среды выполнения Go. Эксперимент сравнивает две реализации моделей при
  общем для них планировщике языка.

\subsection{Программные платформы эксперимента}
\label{subsec:software-platforms}

Программа выполняется под Ubuntu на системе x86 и под macOS 26 на MacBook Pro с M1 Pro.
  Для каждой платформы используется нативный исполняемый файл. На Apple silicon это исключает
  трансляцию x86-64-кода через Rosetta~\cite{apple_rosetta}. Состав платформ приведён в
  таблице~\ref{tab:software-platforms}.

% TODO: перед итоговыми измерениями дополнить протокол версиями ОС, ядра Linux и Go, архитектурой
% бинарного файла, значением GOMAXPROCS и режимом питания.

\begin{table}[htbp]
  \centering
  \caption{Программные платформы эксперимента}
  \label{tab:software-platforms}
  \begin{tabularx}{\textwidth}{|p{3.4cm}|X|X|}
    \hline
    Параметр & x86-система & ARM-система \\
    \hline
    Операционная система & Ubuntu Linux & macOS 26 \\
    \hline
    Архитектура программы & x86-64 (Intel 64) & arm64 (AArch64) \\
    \hline
    Язык реализации & Go & Go \\
    \hline
    Режим исполнения & Нативный код x86-64 & Нативный код arm64 без Rosetta \\
    \hline
  \end{tabularx}
\end{table}

Основное сравнение проводится отдельно на каждой машине: для вариантов 1:1 и M:N используются
  одни входные данные, версия Go и параметры запуска. Сопоставление результатов между машинами
  служит дополнительной проверкой. Поскольку вместе с процессором меняются ОС и память,
  различия между x86- и ARM-системами нельзя приписать только архитектуре команд или только ОС.
